
\subsubsection{3.1 Agency Proxy Extraction}

Four agency proxies are defined, motivated by HumanAgencyBench dimensions:

\begin{itemize}
\item \textbf{Clarifying Questions.} Binary indicator for responses containing question structures that seek clarification before providing advice.
\item \textbf{Option Enumeration.} Count of explicitly enumerated alternatives presented to the user.
\item \textbf{Epistemic Hedging.} Ratio of uncertainty markers to total tokens.
\item \textbf{Explicit Deferral.} Binary indicator for responses that explicitly defer to user judgment.
\end{itemize}

For each proxy, a TF-IDF + Logistic Regression classifier is trained using regex-generated labels as ground truth. TF-IDF vectorization uses n-gram range (1,2) with maximum 5000 features. Logistic Regression uses C=1.0.

\subsubsection{3.2 BAI Computation}

The Bidirectional Alignment Index combines proxy probabilities with length normalization:

$$\text{BAI} = \frac{1}{4} \sum_{p \in \text{proxies}} P(p|\text{response}) \cdot \frac{1}{1 + 0.1 \cdot \log(\text{word\_count})}$$

\subsubsection{3.3 Adversarial Probing for Representational Independence}

To test whether BAI occupies a representationally independent subspace, gradient reversal training is employed with a shared encoder, a reward probe trained with standard gradient descent, and a BAI probe trained with gradient reversal.

If BAI occupies an independent subspace, the BAI probe should maintain high performance (AUROC $\geq$ 0.7) even after gradient reversal suppresses reward-predictive variance.

The gradient reversal layer uses DANN sigmoid scheduling with alpha ramping from 0 to 1 over the first epoch. Loss weighting uses $\lambda =1.0$ for the reward term.

\subsubsection{3.4 Disagreement Analysis}

Systematic disagreement is quantified through quartile-based analysis. Responses are assigned to quadrants after z-score standardization: high-BAI/high-reward (HH), high-BAI/low-reward (HL), low-BAI/high-reward (LH), and low-BAI/low-reward (LL). The disagreement rate is computed as (HL + LH) / total.

\subsubsection{3.5 Semantic Coherence Validation}

Disagreement responses are clustered using BERTopic with all-MiniLM-L6-v2 embeddings, UMAP (5 components, 15 neighbors), and HDBSCAN (minimum cluster size 50). Topics are examined for agency vocabulary (clarify, prefer, might, option, consider, perhaps, recommend, suggest, defer, decide).
